% Revised GeoMapBench geospatial task taxonomy (23 leaves)
\documentclass{article}
\usepackage[margin=0.7cm,landscape]{geometry}
\usepackage{forest}
\usepackage{xcolor}
\usepackage{graphicx}
\usepackage{ifthen}
\useforestlibrary{edges}
\pagestyle{empty}

% Uses the user's icons when present and remains compilable without them.
\newcommand{\Ticon}{%
  \IfFileExists{text_vector.png}%
  {\raisebox{-0.40\height}{\includegraphics[height=1.8em]{text_vector.png}}}%
  {\raisebox{-0.15em}{\fbox{\scriptsize T}}}}
\newcommand{\Iicon}{%
  \IfFileExists{image_vector.png}%
  {\raisebox{-0.40\height}{\includegraphics[height=1.8em]{image_vector.png}}}%
  {\raisebox{-0.15em}{\fbox{\scriptsize I}}}}
\newcommand{\TI}{\Ticon\;\Iicon}
\newcommand{\io}[2]{\\{\scriptsize Input: #1 \quad Output: #2}}

\definecolor{hidden-black}{RGB}{0,0,0}

\begin{document}

\begin{figure}[p]
\centering
\resizebox{!}{0.78\textheight}{%
\begin{forest}
  forked edges,
  for tree={
    child anchor=west,
    parent anchor=east,
    grow'=east,
    anchor=west,
    base=left,
    rectangle,
    draw=hidden-black,
    rounded corners,
    minimum height=2em,
    minimum width=4em,
    edge+={gray!55, line width=0.9pt},
    s sep=2.5pt,
    inner xsep=0.4em,
    inner ysep=0.5em,
    line width=0.75pt,
    text width=11em,
    where level=1{fill=yellow!20, font=\normalsize, text width=23em}{},
    where level=2{fill=red!15, font=\small, text width=16em}{},
    where level=3{fill=green!20, font=\small, text width=21em}{},
    where level=4{fill=blue!12, font=\footnotesize, text width=18em}{},
  },
  [GeoMapBench Comprehensive Geospatial Task Taxonomy
    [1. Visual Perception and Interpretation
      [Map and Imagery Understanding
        [Cartographic Symbol Recognition\io{\Iicon}{\Ticon}
          [Identify the point-of-interest meaning of a cartographic symbol.]
        ]
        [Map Text Detection{,} Recognition{,} and Grouping\io{\Iicon}{\TI}
          [Detect word polygons{,} transcribe text{,} and group words into map labels.]
        ]
        [Map Label-to-Feature Anchoring\io{\TI}{\Ticon}
          [Associate a visible map label with the correct point{,} line{,} or polygon feature.]
        ]
        [Remote-Sensing Scene Classification\io{\TI}{\Ticon}
          [Assign a scene-level land-use or land-cover class to an imagery patch.]
        ]
        [Object Presence and Counting\io{\TI}{\Ticon}
          [Determine whether a geographic object is present and count visible instances.]
        ]
      ]
      [Geospatial Feature Delineation
        [Dense Land-Cover Labeling\io{\TI}{\Iicon}
          [Assign a semantic land-cover class to every pixel or patch.]
        ]
      ]
      [Spatiotemporal Visual Comparison
        [Change Localization\io{\TI}{\TI}
          [Localize object-level or pixel-level change between aligned observations.]
        ]
        [Temporal Scene Matching\io{\TI}{\Ticon}
          [Decide whether two observations show the same geographic area across time.]
        ]
      ]
      [Image-to-Location Grounding
        [Visual Geolocation\io{\Iicon}{\Ticon}
          [Infer city{,} country{,} and approximate coordinates from visual geographic cues.]
        ]
      ]
    ]
    [2. Spatial Reasoning and Geometric Processing
      [Coordinate and Reference-Frame Handling
        [Coordinate Transformation\io{\Ticon}{\Ticon}
          [Convert coordinates between declared coordinate reference systems.]
        ]
      ]
      [Spatial Relationship Reasoning
        [Metric Distance Computation\io{\TI}{\Ticon}
          [Compute geodesic or geometry-aware distances between spatial entities.]
        ]
        [Topological and Directional Reasoning\io{\TI}{\Ticon}
          [Infer within{,} contains{,} touches{,} and relative cardinal direction relations.]
        ]
      ]
      [Network and Route Analysis
        [Spatial Graph Construction\io{\TI}{\Ticon}
          [Convert road or path centerlines into a connected{,} weighted routable graph.]
        ]
        [Shortest-Path Optimization\io{\TI}{\TI}
          [Return a valid least-cost path under a declared edge-cost function.]
        ]
        [Isochrone / Service-Area Computation\io{\TI}{\TI}
          [Return the area reachable from an origin within a time or distance budget.]
        ]
      ]
    ]
    [3. Semantic Interpretation and Analytical Reasoning
      [Geo-Semantics and Grounding
        [Toponym Recognition\io{\Ticon}{\Ticon}
          [Detect place-name mentions and their exact spans in natural-language text.]
        ]
        [Geo-Entity Typing\io{\Ticon}{\Ticon}
          [Classify a geographic mention as a populated place{,} water body{,} terrain{,} route{,} etc.]
        ]
        [Textual Spatial-Relation Extraction\io{\Ticon}{\Ticon}
          [Extract explicit or inferred relations such as north of{,} inside{,} near{,} or touching.]
        ]
      ]
      [Comparative and Explanatory Geospatial Reasoning
        [Cross-Entity Comparison\io{\Ticon}{\Ticon}
          [Compare two locations or regions using a declared metric{,} year{,} and unit.]
        ]
      ]
    ]
    [4. Spatial Intelligence Tools Suite
      [Socioeconomic and Environmental Inference
        [Environmental Layer Identification\io{\Ticon}{\Ticon}
          [Identify a climate{,} biome{,} or ecoregion category at a specified location.]
        ]
        [Population and Density Estimation\io{\Ticon}{\Ticon}
          [Retrieve or aggregate population density for a specified region and year.]
        ]
      ]
      [Geoscience Knowledge Reasoning
        [Geologic / Geomorphic Interpretation\io{\Ticon}{\Ticon}
          [Identify and interpret mapped lithology{,} stratigraphic age{,} or landform information.]
        ]
        [Geographic Fact Reasoning\io{\Ticon}{\Ticon}
          [Answer geographic questions requiring spatial relations and structured world knowledge.]
        ]
      ]
    ]
  ]
\end{forest}%
}
\caption{Revised GeoMapBench taxonomy. Text/coordinates/serialized vectors are represented by T; raster imagery and rendered maps by I. The former map-text leaf is split into measurable text-grouping and feature-anchoring tasks; scene classification, object presence/counting, and visual geolocation are added.}
\label{fig:geomapbench_taxonomy_revised}
\end{figure}

\end{document}
